\documentclass[12pt,a4paper]{article}

% ==========================================
% PACOTES
% ==========================================

\usepackage[utf8]{inputenc}
\usepackage[T1]{fontenc}
\usepackage[brazil]{babel}
\usepackage{amsmath}
\usepackage{geometry}
\usepackage{newtxtext,newtxmath}
\usepackage{setspace}
\usepackage[
    colorlinks=true,
    linkcolor=black,
    urlcolor=blue
]{hyperref}

\geometry{a4paper, margin=2cm}
\onehalfspacing

% ==========================================
% DOCUMENTO
% ==========================================

\begin{document}
\pagenumbering{gobble}

% ==========================================
% CAPA
% ==========================================

\begin{titlepage}
    \centering

    {\large Universidade Federal de São João del-Rei -- UFSJ}\\
    {\large Departamento de Ciência da Computação -- DCOMP}\\
    {\large Computação Paralela}\\

    \vspace*{5cm}

    {\Large
        \textbf{\MakeUppercase{Paralelização do Crivo de Eratóstenes com MPI}}
    }\\[2cm]

    {\normalsize
        Autor:\\[0.3cm]
        Samuel de Souza Carvalho Almeida
    }

    \vfill

    {\normalsize Professor: Rafael Sachetto Oliveira}\\[1cm]
    {\large São João del-Rei -- MG}\\
    {\large 2026}
\end{titlepage}

\pagenumbering{roman}

% ==========================================
% SUMÁRIO
% ==========================================

\tableofcontents
\newpage
\pagenumbering{arabic}

% ==========================================
% 1. INTRODUÇÃO
% ==========================================

\section{Introdução}

Este trabalho apresenta a implementação sequencial e a paralelização do
algoritmo Crivo de Eratóstenes para identificar os números primos existentes
no intervalo de 2 até um limite inteiro $N$. A versão paralela utiliza MPI e o
modelo mestre--escravo, no qual o processo mestre distribui o trabalho,
coordena a comunicação e reúne os resultados produzidos pelos processos
escravos.

O desenvolvimento será acompanhado pela análise do desempenho sequencial, pela
identificação das partes de maior custo computacional e pela avaliação dos
ganhos obtidos com a paralelização. Serão considerados diferentes valores de
$N$ e diferentes quantidades de processos.

% ==========================================
% 2. ALGORITMO SEQUENCIAL
% ==========================================

\section{Algoritmo sequencial}

O Crivo de Eratóstenes começa com uma lista que representa os números de 2 até
$N$. A partir do primeiro número primo, inicialmente 2, seus múltiplos são
marcados como compostos. O próximo número ainda não marcado é considerado
primo e o mesmo procedimento é repetido.

Somente os candidatos a primo menores ou iguais a $\sqrt{N}$ precisam ser
avaliados. Entretanto, seus múltiplos são marcados em todo o intervalo até
$N$. Ao final, todos os números que não foram marcados são primos.

Na implementação sequencial, o vetor de marcações é alocado com
\texttt{calloc}, de modo que todas as posições começam indicando que o número
ainda não foi identificado como composto. A marcação dos múltiplos de um primo
$p$ começa em $p^2$, pois os múltiplos menores já foram processados por primos
anteriores.

% ==========================================
% 3. PERFIL SEQUENCIAL
% ==========================================

\section{Perfil de desempenho sequencial}

O perfil da versão sequencial será obtido com a ferramenta \texttt{gprof}. A
análise deverá identificar o tempo gasto na criação do crivo, na marcação dos
números compostos e na contagem dos primos encontrados. Esses resultados serão
utilizados como referência para a escolha dos trechos paralelizáveis.

% ==========================================
% 4. OPORTUNIDADES DE PARALELIZAÇÃO
% ==========================================

\section{Oportunidades de paralelização}

A principal oportunidade de paralelização está na marcação dos números
compostos. Depois que os primos até $\sqrt{N}$ são conhecidos, diferentes
partes do intervalo podem ser processadas de maneira independente. Assim, o
intervalo $[2,N]$ pode ser dividido em blocos contínuos, com um bloco atribuído
a cada processo escravo.

Cada escravo mantém apenas o vetor correspondente ao seu próprio intervalo.
Essa organização reduz a memória utilizada por processo e evita que processos
diferentes escrevam na mesma região de memória.

% ==========================================
% 5. ESTRATÉGIA PARALELA
% ==========================================

\section{Estratégia de paralelização}

O processo mestre, de rank 0, será responsável por ler $N$, calcular os primos
até $\sqrt{N}$ e dividir o intervalo entre os processos escravos. A lista de
primos-base será enviada a todos os processos por meio de uma operação de
\textit{broadcast}. O mestre será utilizado somente para coordenação, conforme
o modelo descrito no enunciado e apresentado em sala. Dessa forma, a marcação
dos números compostos será realizada pelos $P-1$ processos escravos.

A divisão utilizará blocos estáticos e contínuos. Essa escolha reduz a
quantidade de mensagens e é adequada ao crivo, pois os blocos apresentam custos
de processamento semelhantes. Quando a divisão do intervalo não for exata, o
último escravo receberá também as posições restantes. Essa solução foi adotada
por sua simplicidade, uma vez que o pequeno desequilíbrio não é relevante para
a primeira implementação.

Cada escravo utilizará os primos-base para marcar os compostos existentes em
seu bloco. O mestre informará os limites inicial e final de cada intervalo por
meio de uma mensagem \texttt{MPI\_Send}. O envio explícito dos dois limites
mantém a responsabilidade da divisão no mestre e torna simples uma futura
alteração na política de distribuição.

O vetor de marcações local utilizará valores do tipo \texttt{bool}. Essa
representação foi escolhida por expressar diretamente os dois estados
necessários, primo possível ou composto, mantendo a primeira versão simples e
legível. Uma representação compactada em bits poderá ser avaliada futuramente
caso o consumo de memória se torne relevante.

Para cada primo-base $p$, considerando $a$ como o início do bloco, o primeiro
múltiplo marcado será calculado por

\begin{equation}
    \max\left(p^2, \left\lceil\frac{a}{p}\right\rceil p\right).
\end{equation}

Esse cálculo evita marcar o próprio primo quando ele pertence ao bloco e
também evita repetir o trabalho que já foi realizado por primos menores.

Depois do processamento, cada escravo enviará primeiro a quantidade de primos
encontrados e, em seguida, a lista local. Esse protocolo permite que o mestre
aloque exatamente a memória necessária para receber listas de tamanhos
diferentes. O vetor de marcações não será transmitido, reduzindo a quantidade
de dados enviados e correspondendo diretamente ao conteúdo que será gravado no
arquivo de saída.

Como os blocos representam intervalos contínuos e ordenados, o mestre receberá
e concatenará os resultados pela ordem dos ranks. Assim, a lista final já
estará em ordem crescente e não precisará de uma etapa adicional de ordenação.
O arquivo \texttt{resultados/paralelo/primos.txt} armazenará um número primo
por linha, facilitando sua escrita e leitura posterior.

O tempo da execução paralela será medido no processo mestre, começando antes da
distribuição dos dados e terminando após o recebimento do último resultado. O
tempo utilizado para escrever o arquivo de saída não será incluído nessa
medição. Antes do início do cronômetro, uma barreira garantirá que todos os
processos estejam prontos. A medição utilizará a função \texttt{MPI\_Wtime}.

Execuções com mais escravos do que números no intervalo serão rejeitadas. Essa
validação evita processos sem trabalho e mantém o particionamento simples. Os
erros identificados antes da distribuição serão comunicados de forma
coordenada para que todos os processos finalizem normalmente. Falhas
irrecuperáveis ocorridas durante o processamento, como erros de alocação nos
escravos, utilizarão \texttt{MPI\_Abort} para impedir que os demais processos
permaneçam bloqueados aguardando mensagens.

As versões sequencial e paralela permanecerão em arquivos separados nesta
etapa. Essa organização evita uma modularização prematura enquanto as
responsabilidades compartilhadas entre as duas implementações ainda não estão
completamente definidas.

% ==========================================
% 6. AVALIAÇÃO DE DESEMPENHO
% ==========================================

\section{Avaliação de desempenho}

A avaliação comparará os tempos das versões sequencial e paralela. O
\textit{speedup} será calculado por

\begin{equation}
    S(P) = \frac{T_{sequencial}}{T_{paralelo}(P)},
\end{equation}

em que $P$ representa a quantidade total de processos MPI. A eficiência será
calculada por

\begin{equation}
    E(P) = \frac{S(P)}{P}.
\end{equation}

Os experimentos deverão variar tanto o valor de $N$ quanto a quantidade de
processos, permitindo observar o custo da comunicação e a capacidade de
escalabilidade da solução.

Ao final de cada rodada, o valor de $N$, a quantidade de processos, a
quantidade de primos e o tempo medido serão acrescentados ao arquivo
correspondente. A versão sequencial utilizará
\texttt{resultados/sequencial/execucoes.csv}, enquanto a versão paralela
utilizará \texttt{resultados/paralelo/execucoes.csv}. Os dois arquivos terão as
mesmas colunas, facilitando a comparação posterior. Essa gravação será
realizada depois do encerramento do cronômetro e não fará parte do tempo de
execução registrado.

Os experimentos serão automatizados por um script que executará diferentes
valores de $N$ e diferentes quantidades de processos. Cada configuração será
repetida para reduzir a influência de variações ocasionais no tempo medido. O
script compilará os programas uma única vez e acrescentará todas as rodadas aos
arquivos de resultados, sem apagar medições anteriores.

Serão utilizados 70 valores de $N$, distribuídos entre $10^3$ e $10^9$, e
configurações com 2, 3, 4, 6, 8 e 12 processos MPI. A quantidade de
processadores lógicos disponíveis será detectada antes dos experimentos. As
configurações que excederem esse valor serão mantidas com
\textit{oversubscription}, permitindo observar o comportamento da aplicação
quando há mais processos MPI do que unidades de processamento disponíveis.

Durante a série de desempenho, a gravação de \texttt{primos.txt} será
desabilitada para evitar a escrita repetida de arquivos com milhões de linhas.
Essa opção não altera o processamento medido: os primos ainda serão calculados,
enviados pelos escravos e recebidos pelo mestre. As execuções individuais
continuarão gerando a lista completa normalmente.

A série completa poderá ser executada manualmente pelo GitHub Actions. Ao final
da execução, os arquivos CSV e um relatório em Markdown serão disponibilizados
como artefatos. O relatório registrará o sistema operacional, o processador, a
quantidade de núcleos físicos e lógicos, a memória, o compilador e a versão do
MPI presentes no runner. Essas informações serão consideradas na comparação
dos experimentos, pois o hardware dos runners hospedados pode variar entre
execuções.

% ==========================================
% 7. CONCLUSÕES
% ==========================================

\section{Conclusões}

Esta seção será completada após a implementação da versão paralela e a execução
dos experimentos de desempenho.

\end{document}
